
% Document settings
\documentclass[11pt]{article}
\usepackage[margin=1in]{geometry}
\usepackage{graphicx}
\usepackage{wrapfig}
\usepackage{multirow}
\usepackage{setspace}
\usepackage{hyperref}
\usepackage{xurl}
\usepackage{booktabs}
\pagestyle{plain}
\setlength\parindent{0pt}

\begin{document}

\Huge INF632 In Class Discussion
\normalsize ~

\noindent\rule{\textwidth}{1pt}\vspace{.5em}

\LARGE Trends in Wearables for Health
\normalsize ~

In this week's readings, We focused on trends in wearable technologies for health.  Several challenges were mentioned.  I'd like you to start with discussing the challenges mentioned in ``Wearable Technology in the Management of Chronic Diseases: A Growing Concern.''  Ask yourself the following for each of these four areas.
\begin{itemize}
		\item Have you experienced anything related to this challenge before?
		\item Do you see this as specific to any devices you are familiar with?
		\item How addressable do you think this challenges might be?
		\item Why do you think this is a yet unsolved problem?
	\end{itemize}

\begin{enumerate}
\item Device Accuracy
\vspace{0.5in}
\item User Engagement
\vspace{0.5in}
\item Data Interpretation
\vspace{0.5in}
\item Data Standardization
\vspace{0.5in}
\item Data Privacy
\vspace{0.5in}

\end{enumerate}

In the context of ``Wearable sensors for health monitoring: Current applications, trends, and future directions,'' ...

What intersection do you see with the earlier concerns and the two different major classes as outlined in this paper?  As a reminder, those were {\bf wearable biochemical biosensors} and {\bf wearable physical sensors}.

\vspace{0.5in}

\newpage
This paper provided a comprehensive background on the following technologies.  Which of these do you see as having the most potential for adoption, and why?

\begin{itemize}
	\item Tear-based wearable sensors (\S2.1)
	\item Sweat-based wearable sensors (\S2.2)
	\item Saliva-based wearable sensors (\S2.3)
	\item Epidermal sensors (\S2.4)
\end{itemize}
\vspace{0.5in}

Now let's contrast these with non-biological physical sensors (\S2.5).  Do you see biochemical sensors replacing physical sensors?  Augmenting?  Or otherwise?  Why?  {\em What am I really getting at here? Consider the future of wearables from your own perspective, as a prospective user.}

\vspace{0.5in}

Now in the context of ``Advances in Wearable Biosensors for Healthcare: Current Trends, Applications, and Future Perspectives'' I want you to consider the role of machine learning and artificial intelligence in processing the voluminous data provided by these sensors.

Many current generative AI models are typically ``generalize'' with an option for personalized further learning.  Does this approach work with health data?  Why, or why not?  Remember the areas of concern from above...

\vspace{0.5in}

In \S4.3 of this paper, the authors discuss wellness and lifestyle with regard to wearables.  This is health, but from a non-clinical perspective.  Consider the areas of concern from above in this context.  Do these change?  Specifically, consider adoption and data privacy - does using a wearable for wellness, not clinical diagnosis, change your perspective on how you would use such a wearable or what you feel would be an acceptable use of your data?

\begin{itemize}
	\item Device Accuracy
	\item User Engagement
	\item Data Interpretation
	\item Data Standardization
\end{itemize}
\vspace{0.5in}

Finally, what ``killer application'' of wearables would be so motivating to you that all concerns thus far would seem moot point, or maybe less relevant enough that you could over look them?  Why?  \S5.3 starts to get at this.  Let's also ask the opposite question.  Is there any application that raises the bar of concerns?

\end{document}
